\loai{Áp dụng công thức}
\begin{vidu} %[2P4-2.4B][Loại1]
Người ta muốn tạo ra từ trường có cảm ứng từ $B = 250.10^{-5}$ T bên trong một ống dây, mà dòng điện chạy trong mỗi vòng của ống dây chỉ là 2 A thì số vòng quấn trên ống phải là bao nhiêu? Biết ống dây dài 50 cm.
\choice
{7490 vòng}
{4790 vòng}
{479 vòng}
{\True 497 vòng}
\loigiai{Cảm ứng từ trong lòng ống dây: $B=4\pi.10^{-7}nI=4\pi.10^{-7}.\dfrac{N}{l}.I\Rightarrow N=\dfrac{Bl}{4\pi.10^{-7}.2}=497$ vòng}
\end{vidu}
\loai{Mối quan hệ giữa cảm ứng từ với cường độ dòng điện}
\begin{vidu} %[2P4-2.4B][Loại2]
Một ống dây có dòng điện 10 A chạy qua thì cảm ứng từ trong lòng ống là 0,2 T. Nếu dòng điện trong ống là 20 A thì độ lớn cảm ứng từ trong lòng ống là
\choice
{\True 0,4 T}
{0,8 T}
{1,2 T}
{0,1 T}
\loigiai{Ta có: $B=4\pi.10^{-7}.\dfrac{N}{R}I\Rightarrow \dfrac{B_2}{B_1}=\dfrac{I_2}{I_1}\Leftrightarrow \dfrac{20}{10}=\dfrac{B_2}{0,2}\Rightarrow B_2=0,4\ T$}
\end{vidu}
\loai{Xác định số vòng dây theo cách cuốn}
\begin{vidu} %[2P4-2.4B][Loại3]
Một ống dây được cuốn bằng loại dây mà tiết diện có bán kính 0,5 mm sao cho các vòng sát nhau. Khi có dòng điện 20 A chạy qua thì độ lớn cảm ứng từ trong lòng ống dây là
\choice
{4 mT}
{8 mT}
{\True $8\pi\ mT$}
{$4\pi\ mT$}
\loigiai{\begin{itemize}
\item Số vòng dây trên một mét là: $N=\dfrac{1}{2R}=1000$ vòng.
\item Độ lớn cảm ứng từ trong lòng ống dây: $B=4\pi.10^{-7}NI=8\pi.10^{-3}\ T$.
\end{itemize}}
\end{vidu}
\begin{vidu} %[2P4-2.4K][Loại3]
Một ống dây thẳng dài 20 cm, đường kính 2 cm. Một dây dẫn có vỏ bọc cách điện dài 3,14 m (tiết diện dây dẫn rất nhỏ) được quấn đều theo chiều dài ống. Ống dây không có lõi và đặt trong không khí. Lấy $\pi = 3,14$. Cường độ dòng điện đi qua dây dẫn là 0,5 A. Cảm ứng từ trong ống dây là
\choice
{$6,28.10^{-3}$ T}
{\True $1,57.10^{-4}$ T}
{$1,57.10^{-3}$ T}
{$6,28.10^{-3}$ T}
\loigiai{\begin{itemize}
\item Chu vi 1 vòng dây được quấn là: $P = \pi D = 6,28\ cm$.
\item Số vòng dây được quấn trên ống là: $N=\dfrac{L}{P}=50$ vòng.
\item $B=4\pi .10^{-7}\dfrac{N}{\ell}I=1,57 .10^{-4}$ T.
\end{itemize}}
\end{vidu}
\loai{Xác định Hiệu điện thế hai đầu ống}
\begin{vidu} %[2P4-2.4K][Loại4]
Một sợi dây đồng có đường kính 0,4 mm, điện trở $R = 1,1\ \Omega$ , lớp sơn cách điện bên ngoài rất mỏng. Dùng sợi dây này để quấn một ống dây. Các vòng được quấn sát nhau. Đặt điện áp U không đổi vào hai đầu ống dây thì cảm ứng từ trong ống dây là $6,28.10^{-3}$ T. Lấy $\pi = 3,14$. Điện áp U có độ lớn là
\choice
{4,4 V}
{\True 2,2 V}
{3,3 V}
{1,1 V}
\loigiai{Ta có: $B=4. \pi .10^{-7}\dfrac{I}{d}\Rightarrow 6,28.10^{-3}=4. 3,14.10^{-7}\dfrac{I}{0,4.10^{-3}}\Rightarrow I=2\ A \Rightarrow U = IR = 2,2$ V}
\end{vidu}
\loai{Xác định suất điện động của nguồn}
\begin{vidu} %[2P4-2.4K][Loại5] 
\impicinpar{Cho mạch điện có sơ đồ như hình bên: L là một ống dây dẫn hình trụ dài 10 cm, gồm 1000 vòng dây, không có lõi, được đặt trong không khí; điện trở $R=5\ \Omega $; nguồn điện có suất điện động $\mathcal{E}$ và điện trở trong $r=1\ \Omega $. Biết đường kính của mỗi vòng dây rất nhỏ so với chiều dài của ống dây. Bỏ qua điện trở của ống dây và dây nối. Khi dòng điện trong mạch ổn định thì cảm ứng từ trong ống dây có độ lớn là $2,51.10^{-2}\ T$. Giá trị của $\mathcal{E}$ là}{\begin{tikzpicture}[scale=0.7,thick,>=stealth']
\draw[snake=coil,segment length=3pt,segment amplitude=10pt,line before snake=2pt,line after snake=2pt] (0,0) -- (5,0);
\draw(0,0)--(-0.5,0)--(-0.5,-2)--(2.4,-2)(2.4,-1.6)--(2.4,-2.4)(2.6,-1.8)--(2.6,-2.2)node[below, xshift=-3pt]{$\mathcal{E},\ r$}(2.6,-2)--(5.5,-2)--(5.5,0)--(5,0);
\end{tikzpicture}}
\choice
{8 V}
{24 V}
{6 V}
{\True 12 V}

\loigiai{\begin{itemize}
\item Cảm ứng từ trong lòng ống dây khi có dòng điện I chạy qua được xác định bởi biểu thức: $B=4\pi.10^{-7}\dfrac{NI}{\ell}\Rightarrow I=\dfrac{2,51.10^{-2}.0,1}{4\pi 10^{-7}.1000}=2\ A$
\item Suất điện động của nguồn $\mathcal{E} =I(R+r)=2.(5+1)=12\ V$.
\end{itemize}}
\end{vidu}